\section{Introduction}

La remontée d'une bulle de gaz dans un liquide au repos est l'un des problèmes les
plus anciens de la mécanique des fluides diphasiques. Nous etudions le cas d'une bulle 
remontant dans un écoulement porteur \emph{turbulent}. Les écoulements à bulles
turbulents interviennent dans un grand nombre de configurations naturelles et
industrielles : colonnes à bulles et réacteurs chimiques (aération, fermentation,
procédés d'oxydation), génie des procédés pétroliers, mais aussi processus
environnementaux à grande échelle comme les échanges gazeux à l'interface
océan-atmosphère sous l'effet du déferlement des vagues. Dans toutes ces
situations, la turbulence ambiante modifie la vitesse à laquelle le gaz remonte, 
et donc les temps de séjour, les taux de transfert de masse et, in fine, le rendement
des procédés concernés.

La question posée dans ce travail est simple à énoncer et difficile à trancher :
\emph{la turbulence ambiante accélère-t-elle ou ralentit-elle la vitesse
d'ascension d'une bulle, par rapport à sa vitesse terminale en liquide au repos ?}
Deux intuitions contradictoires coexistent dans la littérature. D'une part, la
turbulence pourrait favoriser la remontée en réduisant transitoirement la traînée
(par analogie avec la réduction de traînée observée sur des sphères solides en
écoulement turbulent) ou en permettant à la bulle d'exploiter des structures
tourbillonnaires porteuses. D'autre part, elle pourrait au contraire freiner la
bulle : en la faisant dévier de sa trajectoire rectiligne, en la faisant
échantillonner préférentiellement des zones de fluide descendant, ou en augmentant
sa traînée moyenne du fait de la non-linéarité de la loi de traînée vis-à-vis des
fluctuations de vitesse relative. Trancher entre ces scénarios est rendu difficile
par le couplage intrinsèque entre la bulle et la turbulence qui l'entoure : la
bulle déforme et perturbe localement l'écoulement porteur, son sillage instationnaire
interagit avec les structures turbulentes incidentes, et sa trajectoire peut devenir
oblique ou instable même en l'absence de toute turbulence dès lors que le nombre de
Galilée dépasse un seuil critique. Isoler l'effet net de la turbulence sur la
vitesse de remontée exige donc une simulation où les autres
paramètres du problème (rapports de densité et de viscosité, tension superficielle,
taille de bulle) sont maîtrisés et maintenus fixes pendant que l'intensité
turbulente varie seule.

\subsection{État de l'art}

Le ralentissement de la remontée d'une bulle par la turbulence ambiante est un
problème débattu de longue date. Sur le plan expérimental, Poorte \& Biesheuvel
\cite{poorte2002} ont mesuré, dans une turbulence de grille quasi homogène et
isotrope, une réduction de la vitesse moyenne d'ascension pouvant atteindre
environ $-35\,\%$ par rapport à la valeur en liquide au repos, établissant
empiriquement que l'effet net de la turbulence est bien un ralentissement et non
une accélération. Sur le plan théorique, Spelt \& Biesheuvel \cite{spelt1997} ont
proposé, dans la limite d'une turbulence de faible intensité, une loi de
correction quadratique de la forme $1 - c\,\beta^2$, où $\beta = u'/u_\infty$ est
le rapport entre la fluctuation de vitesse turbulente $u'$ et la vitesse
terminale en milieu quiescent $u_\infty$ : la réduction relative de vitesse croît
comme le carré de l'intensité turbulente relative au voisinage de $\beta = 0$.
À l'opposé, dans le régime de turbulence forte, Ruth \emph{et al.}
\cite{ruth2021} ont montré que la réduction de vitesse suit asymptotiquement une
loi en $0.37/\mathrm{Fr}'$, où $\mathrm{Fr}' = u'/\sqrt{gD}$ est un nombre de
Froude turbulent construit sur la fluctuation de vitesse et le diamètre de bulle
$D$ : au-delà d'un certain régime, c'est la dynamique de grand nombre de Froude
turbulent, et non plus l'intensité turbulente relative $\beta$, qui contrôle
directement la réduction de vitesse.

Ces deux comportements ont été réunis en une courbe  par Liu, Farsoiya, Perrard \& Deike \cite{liu2024}, au moyen de simulations
numériques directes réalisées avec le même solveur (Basilisk) et sur le même type
de configuration (une bulle unique en turbulence homogène isotrope triplement
périodique) que celle utilisée dans le présent travail. Cette étude sera la
référence  du présent stage : elle relie le régime petite
perturbation ($1-c\beta^2$) au régime grand nombre de Froude turbulent
($0.37/\mathrm{Fr}'$), et surtout elle attribue explicitement le mécanisme du
ralentissement à la \emph{traînée non linéaire} : la loi de traînée $C_D(Re)$
n'étant pas linéaire en la vitesse relative, des fluctuations symétriques de
cette vitesse relative induites par la turbulence produisent, en moyenne, une
traînée plus grande que celle évaluée à la vitesse moyenne, d'où un
ralentissement systématique. Ce mécanisme est explicitement distingué de
l'\emph{échantillonnage préférentiel} (biais statistique par lequel une
bulle plus légère que le fluide environnant a tendance à être capturée par les
zones de vorticité et de fluide descendant). Ce mécanisme, selon Liu \emph{et
al.}, domine plutôt pour des bulles de taille inférieure ou de l'ordre de
l'échelle de Kolmogorov (nombre de Stokes $\mathrm{St} \sim 1$), alors que la
traînée non linéaire domine pour des bulles de taille supérieure à cette échelle
($\mathrm{St} \gg 1$), configuration qui est celle du présent travail.

Le paramètre de contrôle central de ce problème a été mis en évidence par la
revue récente de Legendre \& Zenit \cite{legendre2025} (\emph{Reviews of Modern
Physics}) : c'est le rapport $\beta = u'/u_\infty$ entre la fluctuation de
vitesse turbulente et la vitesse terminale en milieu au repos qui organise
l'ensemble des résultats de la littérature. D'autres travaux de simulation numérique
directe éclairent des aspects connexes du problème : Loisy \& Naso
\cite{loisy2017} ont étudié le cas d'une grosse bulle flottante en turbulence
homogène isotrope, dans un régime de tailles comparable à celui considéré ici ;
Cano-Lozano \& Magnaudet \cite{canolozano2016} ont caractérisé l'instabilité de
trajectoire d'une bulle isolée en liquide au repos, montrant qu'elle devient
oblique ou instationnaire au-delà d'un nombre de Galilée critique
$\mathrm{Ga}_c \approx 44$ ; enfin Rivière \emph{et al.} \cite{riviere2021} ont
caractérisé la fragmentation d'une bulle en turbulence, avec un seuil critique en
nombre de Weber turbulent $\mathrm{We}_c \approx 3$, qui borne le domaine des
intensités turbulentes explorables sans rupture de la bulle.

De cette littérature se dégagent trois éléments structurants pour la présente
étude. Premièrement, le paramètre de contrôle pertinent est $\beta = u'/u_\infty$. Deuxièmement, la réduction relative de vitesse
$u_{\infty,\mathrm{turb}}/u_\infty$ suit deux régimes asymptotiques distincts : une
loi quadratique $1-c\beta^2$ à petit $\beta$, puis une loi en $1/\mathrm{Fr}'$ à
grand nombre de Froude turbulent, les deux étant reliées par la courbe de référence
de Liu \emph{et al.} \cite{liu2024}. Troisièmement, le mécanisme physique
responsable du ralentissement fait débat entre échantillonnage préférentiel et
traînée non linéaire, les deux mécanismes n'étant pas mutuellement exclusifs mais
opérant préférentiellement dans des régimes de taille de bulle différents
(sous-Kolmogorov contre super-Kolmogorov).

%%% Ajouter ici une figure de la courbe maitresse de liu et al

\subsection{Objectifs et plan du rapport}

L'objectif de ce travail est de quantifier, par simulation numérique, la
réduction relative de la vitesse d'ascension $u_{\infty,\mathrm{turb}}/u_\infty$
d'une bulle unique en fonction de l'intensité turbulente relative $\beta$, à
nombre de Bond $\mathrm{Bo}=1$ et nombre de Galilée $\mathrm{Ga}=70$ fixés (donc
au-delà du seuil d'instabilité de trajectoire $\mathrm{Ga}_c\approx44$ identifié
par Cano-Lozano \& Magnaudet \cite{canolozano2016}), et d'identifier lequel des
deux mécanismes candidats 1. échantillonnage préférentiel ou 2. traînée non linéaire domine dans le régime de taille de bulle considéré, en faisant la confrontation aux résultats de Liu, Farsoiya, Perrard \& Deike
\cite{liu2024}.

La suite du rapport est organisée comme suit. La section~2 pose le cadre physique
et adimensionnel du problème (Vaschy-Buckingham, définition de
$\beta$ et de $\mathrm{Fr}'$). La section~3 décrit les méthodes numériques :
solveur Basilisk, méthode VOF conservant la quantité de mouvement, raffinement de
maillage adaptatif par octree, protocole en deux étapes précurseur turbulent puis
injection de la bulle. La section~4 présente la validation du modèle numérique
(convergence en maillage, courants parasites, biais liés à la périodicité du
domaine). La section~5 présente les résultats : effet mesuré de la turbulence sur
la vitesse d'ascension en fonction de $\beta$. La section~6 discute le mécanisme
physique du ralentissement observé. La section~7 conclut et dégage les
perspectives de ce travail.
